%------------------------------------------------------------------------------%
% DISTRIBUTION THEORY                                                          %
%------------------------------------------------------------------------------%
% File  : Distribution_Theory.tex                                              %
% Author: Klaus Hoeflinger, Beethovenstrasse 11, D-73117 Wangen                %
% Email : khoeflinger@t-online.de                                              %
%------------------------------------------------------------------------------%

%------------------------------------------------------------------------------%
%                       DOCUMENT CLASS and INCLUDE FILES                       %
%------------------------------------------------------------------------------%
\documentclass[11pt,twoside]{article}
\usepackage{geometry}
\geometry{paperheight=241mm,
          paperwidth=152mm,
          top=22mm,
          bottom=17mm,
          left=20mm,
          right=20mm,
          textwidth=112mm,
          textheight=202mm,
          headheight=10mm,
          headsep=3mm,
          footskip=9mm,
          includehead,
          includefoot,
          marginparsep=0mm,
          marginparwidth=0mm}

\renewcommand{\baselinestretch}{0.945}\normalsize

\AtBeginDocument{\setlength\abovedisplayskip{2mm}}
\AtBeginDocument{\setlength\belowdisplayskip{2mm}}

%------------------------------------------------------------------------------%
% define title formats                                                         %
%------------------------------------------------------------------------------%
\usepackage{titlesec}

\renewcommand{\thesection}{\arabic{section}}
\renewcommand{\sfdefault}{FiraSans}
\renewcommand{\ttdefault}{pcr}
\renewcommand{\rmdefault}{antiqua}

\titleformat{\part}[display]
{\normalfont \large \rmfamily \bfseries \centering}
{\small{\thepart.}}{2pt}{}

\titleformat{\section}[runin]
{\normalfont \rmfamily \bfseries}
{\thesection}{1pt}{.\;}[.]

%\titleformat{\section}
%{\normalfont \bfseries \small \filcenter}
%{\thesection.\;}{0mm}{}

\titleformat{\subsection}[runin]
{\normalfont \itshape}
{}{0mm}{}

\titleformat{\subsubsection}[runin]
{\normalfont \itshape}
{}{}{}{}

\titleformat{\paragraph}[runin]
{\normalfont}
{}{}{}{}

\titleformat{\subparagraph}[runin]
{\normalfont}
{}{}{}{}

\titlespacing*{\part}          { 0pt}{ 0pt}{ 8pt}
\titlespacing*{\section}       { 0pt}{ 7pt}{ 4pt}
\titlespacing*{\subsection}    { 0pt}{14pt}{ 6pt}
\titlespacing*{\subsubsection} { 0pt}{ 6pt}{12pt}
\titlespacing*{\paragraph}     { 0pt}{ 6pt}{12pt}
\titlespacing*{\subparagraph}  { 0pt}{ 0pt}{12pt}

%------------------------------------------------------------------------------%
% define enumerate parameters                                                  %
%------------------------------------------------------------------------------%
\usepackage{enumitem}
\setlength{\parindent}{3pt}
\setlength{\parskip}  {0pt}

%\setlist{noitemsep}

\setlist[enumerate]{
         leftmargin=12pt,
         rightmargin=0pt,
         itemsep=1pt,
         itemindent=3pt,
         topsep=3pt,
         labelsep=4pt,
         partopsep=0pt,
         labelwidth=0pt,
         listparindent=0pt,
         parsep=0pt}

\setlist[itemize]{
         leftmargin=12pt,
         rightmargin=0pt,
         itemsep=0pt,
         itemindent=1pt,
         topsep=1pt,
         labelsep=4pt,
         partopsep=1pt,
         labelwidth=0pt,
         listparindent=0pt,
         parsep=0pt}

%------------------------------------------------------------------------------%
% define packages                                                              %
%------------------------------------------------------------------------------%
\usepackage{upref}
\usepackage{amssymb}
\usepackage{slantsc}
\usepackage{amsmath}
\usepackage{amsthm}
\usepackage{thmtools}
\usepackage{amscd}
\usepackage{verbatim}
\usepackage{latexsym}
\usepackage{multicol}
\usepackage{graphicx}
\usepackage{setspace}
\usepackage[ngerman,english]{babel}
\usepackage[protrusion=true,expansion=true]{microtype}
\usepackage{tikz}
\usetikzlibrary{arrows,arrows.meta,decorations.markings}
\usepackage{mathrsfs}
\usepackage{amsfonts}
\usepackage{datetime2}
\usepackage{textgreek}
\usepackage{hyperref}
\usepackage{pifont}
\usepackage{relsize}
%\usepackage{biblatex}
\relsize{-0.05}
\usepackage[skip=0mm]{parskip}
\usetikzlibrary{decorations.pathmorphing}

\usepackage{mlmodern}
\usepackage[T5,T1]{fontenc}
\usepackage{ragged2e}

%\usepackage{mathabx}
%\usepackage{times}

%\usepackage{fontspec}
%\setmainfont{Nimbus Roman No9 L}
%\usepackage[T1]{fontenc}

%------------------------------------------------------------------------------%
% define page layout                                                           %
%------------------------------------------------------------------------------%
\usepackage{fancyhdr}
\pagestyle{fancy}
\setlength{\headheight}{30.0pt}
\renewcommand{\headrulewidth}{0pt}

\AtBeginDocument{
  \addtolength\abovedisplayskip{-0.0\baselineskip}
  \addtolength\belowdisplayskip{-0.0\baselineskip}
}

\fancyhead{}
\fancyfoot{}
\addtolength{\headwidth}{0.02\marginparwidth}

\makeatletter
\let\ps@plain\ps@fancy
\makeatother

\renewcommand\sectionmark[1]
{
 \def\sectionname{#1}
 \markright{\thesection #1}
}

\makeatletter
\@addtoreset{section}{part}
\makeatother

%------------------------------------------------------------------------------%
% define footnote without number                                               %
%------------------------------------------------------------------------------%
\newcommand{\myfootnote}[1]{%
\renewcommand{\thefootnote}{}%
\footnotetext{#1}%
\renewcommand{\thefootnote}{\arabic{footnote}}%
}

%------------------------------------------------------------------------------%
% define numbering in equations                                                %
%------------------------------------------------------------------------------%
\numberwithin{equation}{section}

%------------------------------------------------------------------------------%
% define newtheorem                                                            %
%------------------------------------------------------------------------------%
\declaretheoremstyle[
headfont=\scshape, 
headindent = 0mm, %\parindent,
%postheadhook = {\hspace{0mm}\newline},
%postheadspace = \newline, 
spaceabove = 3mm, 
spacebelow = 3mm,
numberwithin=section,
name=Theorem]{khMATH}
\declaretheorem[style=khMATH]{theorem}

\declaretheoremstyle[
headfont=\bfseries, 
headindent = 0mm, %\parindent,
%postheadhook = {\hspace{0mm}\newline},
%postheadspace = \newline, 
spaceabove = 3mm, 
spacebelow = 3mm,
numberwithin=part,
name=Theorem]{khMATH1}
\declaretheorem[style=khMATH1]{theorem1}

\declaretheoremstyle[
headfont=\scshape, 
headindent = 0mm, %\parindent,
%postheadhook = {\hspace{0mm}\newline},
%postheadspace = \newline, 
spaceabove = 3mm, 
spacebelow = 3mm,
numberwithin=section,
name=Corollary]{khMATH1}
\declaretheorem[style=khMATH1]{corollary}

%            name         counter     label              numeration-mode
\theoremstyle{plain}
\newtheorem*{introduction}            {\sc Introduction}
\newtheorem {standard}                {}                 [section]
\newtheorem*{definition}              {\sc Definition}
\newtheorem{satz1}                    {\sc Satz}
\newtheorem{lemma}                    {\sc Lemma}
\newtheorem*{proposition}             {\sc Proposition}
%\newtheorem{corollary}               {\sc Corollary}
\newtheorem*{corollar}                {\sc Korollar}
\newtheorem*{remark}                  {\sc Remark}
\newtheorem*{remarks}                 {\sc Remarks}
\newtheorem*{example}                 {Example}
\newtheorem*{examples}                {Examples}
\newtheorem*{theo}                    {\sc Theorem}

%------------------------------------------------------------------------------%
% define tabular                                                               %
%------------------------------------------------------------------------------%
\usepackage{tabularx}
\newcolumntype{L}[1]{>{\raggedright\arraybackslash}p{#1}}
\newcolumntype{C}[1]{>{\centering\arraybackslash}  p{#1}} 
\newcolumntype{R}[1]{>{\raggedleft\arraybackslash} p{#1}} 

%------------------------------------------------------------------------------%
% define \sh, \zB                                                              %
%------------------------------------------------------------------------------%
\def\dsh{{d.\,h.}}
\def\ie{{i.\,e.}}
\def\zB{z.\,B.}
\renewcommand{\abstractname}{ }

%------------------------------------------------------------------------------%
% change default spacing for cases                                             %
%------------------------------------------------------------------------------%
\makeatletter
\renewcommand*\env@cases[1][1.6]{%
  \let\@ifnextchar\new@ifnextchar
  \left\lbrace
  \def\arraystretch{#1}%
  \array{@{}l@{\quad}l@{}}%
}
\makeatother

\newenvironment{rcases}
  {\left.\begin{aligned}}
  {\end{aligned}\right\rbrace}

%------------------------------------------------------------------------------%
% change default for \predisplaypenalty                                        %
%------------------------------------------------------------------------------%
\predisplaypenalty=990
\allowdisplaybreaks \numberwithin{equation}{section}

%------------------------------------------------------------------------------%
% general notations                                                            %
%------------------------------------------------------------------------------%
\def\*{\star}
\def\={\,=\,}
\def\+{\,+\,}
\def\xsubset{{\,\subset\,}}
\def\xtimes{{\,\times\,}}
\def\xeof{{\,\in\,}}
\def\eq{{\,=\,}}
\def\xto{{\,\to\,}}
\def\eqdef{\,:=\,}
\def\:{{\,:\,}}
\def\ele{{\,\in\,}}
\def\exists{\mbox{ exists }}
\def\and{\mbox{ and }}
\def\for{\mbox{ for }}
\def\fa{\mbox{ for all }}
\def\fe{\mbox{ for each }}
\def\qfa{\quad \mbox{ for all }}
\def\df{\,\dotfill}
\def\iff{if and only if }
\def\wrt{with respect to}

\makeatletter
\newcommand \Df {\leavevmode \cleaders \hb@xt@ .70em{\hss .\hss }\hfill \kern \z@}
\makeatother

\def\sql{\mathfrak{a}}               % sesquilinear form a
\def\DF{B}                           % Dirichlet form a
%------------------------------------------------------------------------------%
% number sets and special elements                                             %
%------------------------------------------------------------------------------%
\def\P{\mathbb{H}}                   % half plane of $\C$
\def\J{I}                            % interval I
\def\N{{\mathbb{N}}}                 % unsigned integers
\def\Q{{\mathbb{Q}}}                 % rational integers
\def\Z{{\mathbb{Z}}}                 % integers
\def\R{{\mathbb{R}}}                 % real numbers
\def\Rn{{\mathbb{R}}^n}              % real numbers in Rn
\def\Rd{{\mathbb{R}}^d}              % real numbers in Rn
\def\C{{\mathbb{C}}}                 % complex numbers
\def\F{{\mathbb{K}}}                 % arbitrary field
\def\Torus{{\mathbb{T}}}             % complex circle line
\def\T{\tau}                         % upper bound of interval (0,t)
\def\sign#1{\operatorname{sign}(#1)} % sign of a number

%------------------------------------------------------------------------------%
% set theory                                                                   %
%------------------------------------------------------------------------------%
\def\G{G}                            % domain $G$
\def\clG{\overline{\G}}              % closure of $G$
\def\bndG{\partial \G}               % boundary of $G$
\def\setA{\mathcal{A}}               % set A
\def\setB{\mathcal{B}}               % set B
\def\setC{\mathcal{C}}               % set C
\def\setM{M}                         % set M
\def\setN{N}                         % set N
\def\setS{\mathcal{S}}               % set S
\def\famF{\mathfrak{F}}              % family of sets
\def\indI{\mathcal{I}}               % index set I
\def\pot#1{\mathfrak{P}(#1)}         % power set

\def\relR{\mathbf{R}}                % relation R
\def\mapf{f}                         % mapping f
\def\mapg{g}                         % mapping g

\def\set#1#2{\{ \, #1 \,: \, #2 \, \}}
\def\tensor#1#2{#1 \otimes #2}
\def\Tens{\oplus}

\def\cal#1{{\mathcal{#1}}}
\def\aut#1{\operatorname{Aut}(#1)}

\def\cross#1#2{#1 {\,\times\,} #2}
\def\multicross#1#2{#1 \times  \ldots \times  #2}
\def\multitens#1#2 {#1 \otimes \ldots \otimes #2}

%------------------------------------------------------------------------------%
% group theory                                                                 %
%------------------------------------------------------------------------------%
\def\1{{\mathsf{1}}}                % unit element of a monoid
\def\0{{\mathsf{0}}}                % unit element of an Abelian monoid
\def\kernel#1{\mathfrak{N}(#1)}     % kernel of a homomorphism
\def\cokernel#1{\mathfrak{C}{(#1)}} % cokernel of a homomorphism
\def\Hom#1{\operatorname{Hom}(#1)}

%------------------------------------------------------------------------------%
% linear algebra                                                               %
%------------------------------------------------------------------------------%
\def\vecV{\mathit{V}}               % vector space V
\def\vecH{\mathit{H}}               % vector space H
\def\vecW{\mathit{W}}               % vector space W
\def\vecU{\mathit{U}}               % subspace U
\def\convC{\mathcal{C}}             % convex set C
\def\basH{\mathfrak{B}}             % Hamel base B

\def\LO#1{\mathscr{L}(#1) }         % algebra of linear transformations

\def\scalar#1#2{ \langle #1,#2 \rangle }
\def\trace#1{\operatorname{tr}(#1)}
\def\dim#1{{\operatorname{dim} \, #1 }}
\def\codim#1{{\operatorname{codim} \, #1 }}
\def\dim{\operatorname{dim}}
\def\codim{\operatorname{codim}}
\def\ind{\operatorname{ind}}

\def\genG{\cal{G}}

\def\algA{\cal{A}}
\def\algB{\cal{B}}
\def\algU{\cal{U}}

\def\idI{\cal{I}}

%------------------------------------------------------------------------------%
% topology                                                                     %
%------------------------------------------------------------------------------%
\def\compK{{K}}          % compact set C
\def\openO{\mathcal{O}}  % open set O
\def\openU{\mathcal{U}}  % open set U
\def\U{\mathcal{U}}      % neighbourhood U
\def\V{\mathcal{V}}      % neighbourhood V
\def\d{\mathit{d}}       % metric function d

\def\interior#1{\mathring{#1}}
\def\closure#1{{\overline{#1}}}

\def\topT{\mathcal{T}}   % topology T
\def\topX{X}             % topological space X
\def\topY{Y}             % topological space Y
\def\topZ{Z}             % topological space Z
\def\metrS{M}            % metric space M

\def\grpG{G}             % topological group G
\def\grpA{A}             % topological group A
\def\grpB{B}             % topological group B
\def\grpC{C}             % topological group C
\def\grpH{H}             % topological group H
\def\grpU{U}             % topological group U
\def\grpV{V}             % topological group V
\def\topG{G}             % topological group G
\def\topH{H}             % topological group H
\def\topU{U}             % topological subgroup U

\def\subS{\cal{S}}
\def\riR{R}              % ring R

%------------------------------------------------------------------------------%
% calculus                                                                     %
%------------------------------------------------------------------------------%
\def\supp#1{\operatorname{supp}(#1)}
\def\singsupp#1{\operatorname{sing \, supp}(#1)}

\def\re{\operatorname{Re}}
\def\im{\operatorname{Im}}
\def\grad{\operatorname{grad}}

%\def\abs#1 {\left|#1\right|}
\def\abs#1 {\lvert #1 \rvert}
%\def\abs#1 {\left|\,#1\,\right|}

%------------------------------------------------------------------------------%
% measure theory                                                               %
%------------------------------------------------------------------------------%
\def\salgA{\Sigma}               % sigma-algebra S
\def\salgB{\cal{B}}              % sigma-algebra B
\def\Borel#1{\cal{B}_{#1}}       % Borel-algebra 
\def\LB{\lambda}                 % Lebesgue-Borel measure

%------------------------------------------------------------------------------%
% function spaces                                                              %
%------------------------------------------------------------------------------%
\def\Lp{L^p(\G)}                 % spaces L_p
\def\Lq{L^q(\G)}                 % spaces L_q
\def\Lh{L^2(\G)}                 % spaces L_2
\def\Li{L^{\infty}(\G)}          % spaces L_infty
\def\Lc{L_{loc}^1(\G)}           % spaces L_1^{loc}
\def\Ck{C^k(\G)}                 % spaces C_k

\def\BV#1{BV(#1) }               % set of functions of bounded variation
\def\AC#1{AC(#1) }               % set of absolutely continuous functions

\def\MP#1{\mathscr{M}^+(#1) }    % set of positive measures
\def\MS#1{\mathscr{M}^-(#1) }    % vector space of finite signed measures
\def\MC#1{\mathscr{M}(#1) }      % vector space of complex measures
\def\MCR#1{\mathscr{M}_{r}(#1) } % vector space of regular complex measures

\def\SobW{W^{k,p}(\G)}                % Sobolev space W
\def\SobO{\mathring{W}^{k,p}(\G)}     % Sobolev space W
%\def\WN#1#2{\mathring{W}^{#1}(#2) }   % Sobolev space W
\def\WN#1#2{W_0^{#1}(#2) }   % Sobolev space W
%\def\HN#1#2{\mathring{H}^{#1}(#2) }   % Sobolev space H
\def\HN#1#2{H_0^{#1}(#2) }   % Sobolev space H
%------------------------------------------------------------------------------%
% functional analysis                                                          %
%------------------------------------------------------------------------------%
\def\snorm#1{\lVert #1 \rVert}
\newcommand{\norm}[1]{\left\lVert #1 \right\rVert}

\def\topV{V}                       % topological vector space V
\def\topW{W}                       % topological vector space W
\def\normS{X}                      % normed space

\def\banX{\mathit{X}}              % Banach space X
\def\banY{\mathit{Y}}              % Banach space Y
\def\banZ{\mathit{Z}}              % Banach space Z
\def\dbanX{\mathit{X^*}}           % dual space of Banach space X
\def\dbanY{\mathit{Y^*}}           % dual space of Banach space Y
\def\dbanZ{\mathit{Z^*}}           % dual space of Banach space Z

\def\banA{\mathit{A}}              % Banach algebra A
\def\banB{\mathit{B}}              % Banach algebra B
\def\banC{\mathit{C}}              % C*-algebra C


\def\domain#1{\mathfrak{D}(#1)}    % domain
\def\range#1{\mathfrak{R}(#1)}     % range

\def\isoJ{\mathfrak{J}}            % isometry J
\def\repA{{\cal{A}_{\sql}}}        % representation operator A
\def\linS{{S}}                     % linear operator S
\def\linG{{G}}                     % linear operator G
\def\linL{{L}}                     % linear operator L
\def\linT{{T}}                     % linear operator A
\def\linA{\mathit{A}}              % linear operator A
\def\linB{{B}}                     % linear operator B
\def\linM{{M}}                     % linear operator M
\def\linU{{U}}                     % linear operator U
\def\linI{{I}}                     % linear operator I
\def\A{\cal{A}}                    % linear differential operator A
\def\BDO{\cal{B}}                  % linear boundary differential operator B
\def\linU{{U}}                     % unitary operator U
\def\linP{{P}}                     % projection P
\def\compA{k}                      % compact operator k
\def\linFR{f}                      % finite rank operator f
\def\fredA{A}                      % Fredholm operator F

\def\HAM{\cal{H}}                  % Hamilton operator
\def\PA{\partial^{\alpha}}         % elementary differential operator
\def\DO{\cal{A}}                   % linear differential operator
\def\BC{\cal{B}}                   % boundary operator
\def\SLO{\cal{L}_{p,q}}            % Sturm-Liouville operator
\def\MO{\cal{L}_{M}}               % momentum operator

\def\HS#1{S(#1)           }        % algebra of Hilbert-Schmidt operators
\def\FR#1{F(#1)           }        % algebra of finite rank operators
\def\TC#1{N(#1)           }        % algebra of trace class operators
\def\BU#1{\mathscr{U}(#1) }        % algebra of unitary linear operators
\def\BO#1{\mathscr{L}(#1) }        % algebra of bounded linear operators
\def\CO#1{\mathscr{C}(#1) }        % algebra of compact linear operators
\def\CL#1{\mathscr{C}(#1) }        % algebra of compact linear operators
\def\FO#1{\Phi(#1) }               % set of Fredholm operators
\def\FK#1#2{\Phi_{#2}(#1) }        % set of Fredholm operators with index k

\def\hilH{\mathit{H}}              % Hilbert space H
\def\clU{\mathit{U}}               % closed subspace U

\def\orthO{\mathfrak{O}}           % orthonormal system O
\def\orthS{\mathfrak{S}}           % orthonormal system S
\def\basB{\mathfrak{B}}            % basis B of a Hilbert space

\def\GL#1 {\mathbf{GL}(#1)}        % linear group
\def\OG#1 {\mathbf{O}(#1)}         % linear group
\def\SOG#1 {\mathbf{SO}(#1)}       % linear group

%------------------------------------------------------------------------------%
% distribution theory                                                          %
%------------------------------------------------------------------------------%
\def\disF{F}                  % distribution F
\def\disG{G}                  % distribution G
\def\disH{H}                  % distribution H
\def\disU{U}                  % distribution U
\def\disT{T}                  % distribution T
\def\SF{C^{\infty}(\G)}       % space of smooth functions
\def\TF{C_c^{\infty}(\G)}     % space of compactly supported smooth functions
\def\TFRd{C_c^{\infty}(\R^d)} % space of test functions
\def\SRd{\mathscr{S}(\R^d)}   % Schwartz space
\def\SR1{\mathscr{S}(\R)}     % Schwartz space
\def\B1#1{C^1(#1) }           % space of continuously differentiable functions
\def\Bm#1{C^m(#1) }           % space of continuously differentiable functions
\def\Ck{C^k(\G)}
\def\TD{\mathscr{S}'(\R^d)}   % space of temperated distributions
\def\E{\cal{E}}               % space of distributions with compact support
\def\D{\mathscr{D}}           % D
\def\FT{\mathcal{F}}          % Fourier transformation F
\def\FT{\mathscr{F}}          % Fourier transformation F

%------------------------------------------------------------------------------%
% other definitions                                                            %
%------------------------------------------------------------------------------%
\def\Proof{\textsc{Proof}}
\def\FRECHET{Fr\'{e}chet}
\def\ARZELA{Arzel\'{a}}
\def\GARDING{G\r{a}rding}
\def\HOELDER{H\"older}
\def\SCHROEDINGER{Schr\"odinger}
\def\JOERGENS{J\"orgens}
\def\WUEST{W\"uest}
\def\KAEHLER{K\"aehler}
\def\HOEFLINGER{H\"oflinger}
\def\POINCARE{Poincar\'{e}}
\def\FA{Functional Analysis}

\def\Author   {K.\,{\HOEFLINGER}}
\def\AuthorUC {K.\,HOEFLINGER}
\def\Version  {1.0}
\def\Email    {khoeflinger@t-online.de}


%\renewcommand\thesection{\Roman{section}}

\def\Title{Distribution Theory}

\titleformat{\section}
{\normalfont \bfseries \filcenter}
{\thesection.\;}{0mm}{}

\titleformat{\subsection}[runin]
{\normalfont \filcenter}
{}{0mm}{}

\titleformat{\subsubsection}[runin]
{\itshape \filcenter}
{}{0mm}{}

\titlespacing*{\section}      {5pt}{12pt}{ 4pt}
\titlespacing*{\subsection}   {0pt}{ 6pt}{12pt}
\titlespacing*{\subsubsection}{0pt}{ 6pt}{ 3pt}

%------------------------------------------------------------------------------%
%                                BEGIN OF DOCUMENT                             %
%------------------------------------------------------------------------------%
\begin{document}
\thispagestyle{empty}

\centerline{\large{\textbf{\Title}}}
\vspace{4pt}
\centerline{\small{{\textsc{\Author}}}}
\vspace{10pt}

\fancyhead[C] {\small{\textsc{\Title}}}
\fancyhead[OL]{\small{\thepage}}
\fancyhead[ER]{\small{\thepage}}

\myfootnote
{
  Version {\Version} from \today;
  Email:  \textit{khoeflinger@\,t-online.de}
}

%------------------------------------------------------------------------------%
%                                 INTRODUCTION                                 %
%------------------------------------------------------------------------------%
Basically,
a \textit{distribution} can be seen as a generalization of a function.
The corresponding subject in {\FA},
called the \textit{theory of distributions},
was developed by \textsc{L.\,Schwartz} in the midth of the 20th century
to describe ideas and concepts of theoretical physics
by an exact mathematical formalism.
Distribution theory can be seen as a direct application
of the theory of locally convex spaces, continuous linear functionals and
transposed linear operators.

%------------------------------------------------------------------------------%
\section{Continuous Linear Operators Revisited}                                %
%------------------------------------------------------------------------------%
As for normed spaces,
the notion of \textit{continuity} is of fundamental interest
for linear operators between topological linear spaces $\vecV$ and $\vecW$.
A linear operator $\linA \: \vecV \xto \vecW$ is continuous
\iff
it is continuous in $0 \xeof \vecV$ or in any other point $x \neq 0$.
But continuity is not equivalent to boundedness in this general setting.

\subsubsection*{} % continuous linear maps between topological linear spaces
If $\vecV$ and $\vecW$ are topological linear spaces,
then $\vecV$ is called \textit{continuously embedded} into $\vecW$,
if there is an injective and continuous linear operator
$\iota \: \vecV \xto \vecW$.
In addition,
if this mapping is surjective and both spaces are {\FRECHET},
then the \textit{open mapping theorem} claims that $\iota^{-1}$ is continuous,
too,
and therefore forms a \textit{linear homeomorphism}
between these spaces (notation: $\vecV \cong \vecW$).

\subsubsection*{} % transposed operators in topological linear spaces
The concept of \textit{transposed operator} is useful to generalize the domains
of linear partial differential operators and Fourier transformation
to the very extensive spaces of distributions.
Let $\vecV$ and $\vecW$ be locally convex and $\vecV^*, \vecW^*$
the dual spaces of them.
Given a continuous linear operator $\linA \: \vecV \xto \vecW$,
each linear functional $f \xeof \vecW^*$ gives rise to
a linear functional $g \xeof \vecV^*$ by the composition $g:= f \circ \linA$.
The assignment $f \mapsto g$ induces a continuous mapping
$\linA^t \, \: \vecW^* \xto \vecV^*$,
called the \textit{trans\-posed operator} of $\linA$.
One says that
the continuous linear functional $g$ is the \textit{pullback} of $f$
along the operator $\linA$.

\subparagraph*{}
The transposed operator $\linA^t$ is also linear.
If $\vecV, \vecW$ are locally convex linear spaces
and $\linA \: \vecV \xto \vecW$ is continuous,
then $\linA^t$ is injective
if and only if $\range{\linA} $ is dense in $\vecW$.
If $\linA$ is bijective,
then $\linA^t$ is bijective, too,
and its inverse $(\linA^t)^{-1}$ is the trans\-posed of $\linA^{-1}$.

\subparagraph*{}
Finally,
if $\vecV$ and $\vecW$ are {\FRECHET} spaces,
then the solvability of the linear equation $\linA x \eq y$ 
is related to certain properties of the transposed operator $\linA^t$:
the linear operator $\linA \: \vecV \xto \vecW$ is surjective
iff
$\linA^t$ is injective and $\range{\linA^t} $ is closed in $\vecV^*$.

%------------------------------------------------------------------------------%
\section{Distributions}                                                        %
%------------------------------------------------------------------------------%
The space $\TF$ of test functions is of fundamental meaning
in distribution theory.
Recall that this space forms a linear subspace of $\SF$
consisting of smooth functions with \textit{compact} support
\[
\supp{\,\phi} \, := \, \overline{\{x \xeof \G: \phi(x) \neq 0\,\} }.
\]
There is neither a countable family of seminorms
making $\TF$ into a {\FRECHET} space nor there is a metric making 
$\TF$ into a F-space.

\subparagraph*{}
But it is still possible to define a locally convex topology $\topT$ on $\TF$
such that this space is complete in the sense of a topological linear space.
Instead of saying how the open sets of this topology $\topT$ look like,
we define $\topT$ by means of a suitable concept of convergence,
known as $\D$-convergence in $\TF$.

\subparagraph*{}
A sequence $(\phi_n)$ in $\TF$ is said to be convergent to $\phi \xeof \TF$,
if there is a compact set $K \xsubset \G$
such that

\begin{itemize}[leftmargin=22mm]
\item[(D-1)]
$\supp{\phi_n} \xsubset K$ for all $n \xeof \N$,

\item[(D-2)]
for each multi-index $\alpha$ the sequence $(\PA \phi_k)$
converges to the function $\PA \phi$ \textit{uniformly} in $K$,
\end{itemize}
the linear space $\TF$ equipped with the topology
induced by $\D$-convergence is usually denoted $\D(\G)$.
We have $\D(\G) \hookrightarrow \E(\G)$,
that is,
the space $\D(\G)$ is continuously embedded into the space $\E(\G)$
{\wrt} their locally convex topologies.

\subsubsection*{} % linear functionals on the space of test functions
The benefit of the function space $\TF$ relies on the fact that,
due to the smoothness and compact support property of its elements,
the integral
\begin{equation}\label{REGULAR_DISTRIBUTION}
\disF_f(\phi) \, := \, \int_{\,\G} f \, \phi \, dx, \quad \quad \phi \xeof \TF
\end{equation}
is well defined for a multiplicity of real or complex-valued functions $f$.
The largest class of functions
for which this integral has a reasonable meaning,
is the class $\Lc$ of \textit{locally integrable functions},
that is,
$f \xeof \Lc$ is Lebesgue-integrable on each compact subset $K \xsubset \G$.
This space incorporates a lot of other well-known function spaces,
for example all the spaces $L^p(\G)$ with $p \geq 1$,
the space $L^{\infty}(\G)$ of essentially bounded functions
and the space $C(\G)$ of continuous functions.
Moreover, 
if $\mu$ is a positive measure on $\G$,
the integral
\begin{equation}\label{RADON_DISTRIBUTION}
\disF_{\mu}(\phi) \, := \, \int_{\,\G} \phi \, d\mu, \quad \, \phi \xeof \TF
\end{equation}
especially makes sense for all the Radon measures $\mu$ on $\G$,
since such measures are finite on compact sets by definition.
Both mappings $\phi \mapsto \disF_f$ and $\phi \mapsto \disF_{\mu}$
as defined in \eqref{REGULAR_DISTRIBUTION} and \eqref{RADON_DISTRIBUTION}
act linearly on $\TF$ and hence give rise to linear functionals.
But both of them are also continuous on the topological linear space $\D(\G)$.

\subsubsection*{} % Schwartz distributions
Let us denote each continuous linear functional $\disF \: \D(\G) \xto \C$ 
a \textit{Schwartz distribution}.
We call $\D'(\G)$ the dual space of $\D(\G)$,
so $\D'(\G)$ is the linear space of Schwartz distributions on the domain $\G$.
One may also define Schwartz distributions {\wrt} $\D$-convergence,
when applied to test functions:
a linear functional $\disF$ on the space $\TF$
is denoted a Schwartz distribution,
if $\disF(\phi_n) \xto 0$ for any $\D$-convergent sequence $(\phi_n)$
converges to the zero function.

\subparagraph*{}
As already mentioned,
all linear functionals derived from locally integrable functions
and Radon measures are continuous and therefore Schwartz distributions.
Indeed,
the space $\D'(\G)$ is extremely extensive
and contains nearly all functions being important in applications,
amongst others the functions
belonging to the well-known spaces $\Lp$, $\Li$, $C(\G)$,
$C(\clG)$ and $C_0(\G)$.

\subparagraph*{}
A distribution is said to be \textit{regular},
if it is the continuous linear functional $\disF_f$
induced by a locally integrable function $f$,
otherwise it is called \textit{singular}.
Well-known examples of singular distributions are
the \textit{Dirac distributions}
\[
\delta_a(\phi) \, := \, \phi(a)
\qfa \phi \xeof \TF
\quad \quad  (a \xeof \G).
\]

\subsubsection*{} % space of distributions with compact support
Regarding the space $\E(\G)$,
which is the space $\SF$ in the locally convex topology
induced by the family of seminorms
\[
\abs{\phi} _{N,n}
\, := \,
\max_{x \xeof K_n} \;
\{ \, \abs{\, (\PA \phi)(x) \,} : \abs{\alpha} \leq N \}
\quad \quad (N,n \xeof \N)
\]
for each $\phi \xeof \SF$,
where $(K_n)$ is a sequence of compact subsets of $\G$ with
\[
K_1 \xsubset K_2 \ldots \xsubset \G,
\quad \quad
\bigcup_{n \xeof \N} \, K_n \= \G
\]
each element $\disF \xeof \E'(\G)$ is called
a \textit{distribution with compact support}.
This notion is founded by the fact
that for such $\disF$ there is a compact set $K_F \xsubset \G$,
called the \textit{support} of $\disF$
such that
$\disF \phi \eq 0$ for all smooth functions $\phi$
having support in $\G \setminus K_F$.

\subparagraph*{}
The denotation for this class of distributions is motivated as follows:
if $\disF_f \xeof \E'(\G)$ is regular,
then the function $f \xeof \Lc$ related to $\disF_f$
has compact support, too.
On the other hand,
each measurable function $f$ with compact support of course
gives rise to a compactly supported distribution $\disF_f \xeof \E'(\G)$.
Since it is always possible to increase a dual space
by decreasing the underlying linear space,
there is a continuous inclusion $\E'(\G) \, \hookrightarrow \, \D'(\G)$,
which is dual to the mentioned embedding
$\D(\G) \, \hookrightarrow \, \E(\G)$.
So we basically may regard
each compactly supported distribution as a Schwartz distribution.

%------------------------------------------------------------------------------%
\section{Convolution and Tensor Products}                                      %
%------------------------------------------------------------------------------%
The usage of distributions for the solution theory of linear partial
differential equations can only be exhausted in full strength
if the notion of \textit{convolution} between elements
of the spaces $\D'(\Rd)$ and $\E'(\Rd)$ has been introduced
as far as possible.

\subparagraph*{}
The convolution $\phi \ast \psi$ of functions $\phi, \psi \xeof \SRd$
is defined according to
\[
(\phi \ast \psi) (x) \, := \, \int_{\,\Rd} \phi(y) \, \psi(x-y) \, dy
\]
and always exists,
hence the space of Schwartz functions forms an algebra
{\wrt} convolution as multiplicative operation.
Furthermore,
convolution between a regular distribution
$\disF_f \xeof \D'(\Rd)$ and a test function $\phi \xeof \D(\Rd)$
is similarly defined:
\[
(\disF_f \ast \phi) (x) \, := \, \int_{\,\Rd} f(y) \, \phi(x-y) \, dy.
\]

\subparagraph*{}
More general,
the convolution $\disF * \disG$ of distributions
$\disF, \disG \xeof \D'(\Rd)$ exists,
if the \textit{support condition} is fulfilled,
meaning that for each compact subset $K \xsubset \Rd$ the set
\[
K_{\disF,\disG} :=
\{(x,y): x \xeof \supp{F}, y \xeof \supp{G}, x+y \xeof K \}
\xsubset \Rd \xtimes \Rd
\]
is compact, too.
In this case we may define the convolution $\disF \ast \disG$ according to
\[
(\disF \ast \disG) \, \phi \, := \,
\iint \limits_{\compK_{\disF,\disG}}
\disF(x)\,\disG(y)\,\phi(x+y) \, dx \, dy, \quad \phi \xeof C_c^{\infty}(\Rd).
\]
The support condition is sufficient
but not necessary for the existence of the convolutional distribution
$\disF \ast \disG$.
Nevertheless,
the condition is always fulfilled
if at least one of the distributions $\disF, \disG$
is compactly supported.
Notice also that
convolution on the space $\D'(\Rd)$ is commutative, if exists.

\subparagraph*{}
Especially the convolution between $\disF \xeof \D'(\G)$
and the Dirac distribution $\delta_0$ always exists and
satisfies the equation $\disF \ast \delta_0 \eq \disF$,
thus $\delta_0$ serves as identity element of this operation.
But convolution is also associative,
hence for distributions $\disF, \disG, \disH \xeof \D'(\Rd)$ the equality
\[
(\disF \ast \disG) \ast \disH \= \disF \ast (\disG \ast \disH)
\]
holds provided that
the involved convolutions involved exist.

\subparagraph*{}
Besides the notion of convolution one also requires
that of \textit{tensor product} between distributions.
Tensor products are useful to transform an inital value
of a distributional differential equation
into an inhomogeneous part of that equation.
Based on the fact
that for $\phi \xeof \D(\R^{d_1})$ and $\psi \xeof \D(\R^{d_2})$ 
the tensor product $\phi \otimes \psi$ defined by
\[
(\phi \otimes \psi)(x,y) \= \phi(x) \cdot \psi(y)\,,
\quad (x,y) \xeof \R^{d_1+d_2}
\]
belongs to the space $\D(\R^{d_1+d_2})$,
there is
a uniquely defined distribution
$
\disF \otimes \disG \xeof \D'(\R^{d_1+d_2})
$
for $\disF \xeof \D'(\R^{d_1})$ and $\disG \xeof \D'(\R^{d_2})$,
which for all test functions
$\phi \xeof \D(\R^{d_1})$ and $\psi \xeof \D(\R^{d_2})$
fulfills the property
\[
\disF \otimes \disG \, (\phi \otimes \psi)
\= \disF(\phi) \cdot \disG(\psi)\,,
\]
also denoted the \textit{tensor product} of $\disF$ and $\disG$.
Performing the tensor product is a commutative and associative operation,
that is
\[
\disF \otimes \disG = \disG \otimes \disF,
\quad \quad
(\disF \otimes \disG) \otimes \disH \=
 \disF \otimes (\disG \otimes \disH).
\]

\subparagraph*{}
Let $\PA$ be the elementary differential expression
given by the multi-index $\alpha$.
Then $\PA$ is a continuous linear operator $\PA \:\D(\G) \xto \D(\G)$,
due to the fact
that the space $\D(\G)$ is endowed with a very fine topology.

\subparagraph*{}
The action $\PA_{\D}$ of the expression $\PA$
on the distributional space $\D'(\G)$
is defined by transposition of the \textit{formally adjoint} expression
$(\PA)^* := (-1)^{\alpha} \PA$,
or equivalently
\begin{equation}\label{DISTRIBUTIONAL_DERIVATIVE}
(\PA_{\D} \disF) \, \phi \, := \, \disF((\PA)^* \phi)
\= (-1)^{\abs{\alpha} } \, \disF(\PA \phi)
\end{equation}
for all $\disF \xeof \D'(\G)$ and $\phi \xeof \D(\G)$.

\subparagraph*{}
It should be emphasized
that this definition is chosen in order
to obtain compatibility with possible classical derivatives
of \textit{regular} distributions.
In fact,
if $F_f \xeof \D'(\G)$ is the distribution generated
by a differentiable function $f \:\G \xto \C$
as defined in \eqref{REGULAR_DISTRIBUTION},
then $\PA_{\D} F_f \eq F_{\PA f}$ holds.

\subparagraph*{}
The mapping $\disF \mapsto \PA_{\D} \disF$ proves to be continuous
on the Schwartz space $\D'(\G)$.
We also have the following compatibility rules
between convolutions and tensor products:
\[
\PA_{\D} (\disF \ast \disG)  \=
(\PA_{\D} \disF) \ast \disG \=
\disF \ast (\PA_{\D} \disG) \,,
\]
\[
 \PA_{\D} (\disF \otimes \disG) =
(\PA_{\D} \disF) \otimes \disG  =
 \disF \otimes (\PA_{\D} \disG).
\]

\subsubsection*{} % Schwartz space
We continue with the introduction of the \textit{Schwartz space} $\SRd$,
the linear space of all smooth functions $\phi \xeof C^{\infty}(\Rd)$
having the additional property
\[
\abs{ \PA \phi(x) \cdot x^{\beta}} \to 0 \,
\for \abs{x} \to \infty
\mbox{ and all multi-indices } \alpha,\beta,
\]
where $\abs{x} $ denotes the Euclidean norm of $x \xeof \Rd$.
A function $\phi \xeof \SRd$ is called \textit{rapidly decreasing}
or briefly a \textit{Schwartz function}.
Hence Schwartz functions and all their derivatives vanish at infinity
faster then the reciprocal of any polynomial.
The most famous example of a Schwartz function is the Hermite function
\[
\cal{H}_0(x) \, := \, e^{- \frac{\abs{x} ^2}{2}}, \quad x \xeof \Rd.
\]

\subparagraph*{}
Rather similar to the space $\SF$,
one can find for the Schwartz space $\SRd$
a family of seminorms as follows:
define for $\phi \xeof C^{\infty}(\Rd)$ and multi-indices $\alpha, \beta$
\begin{equation}\label{X_SRN_FAMILY_OF_SEMINORMS}
p_{\alpha,\beta}(\phi) \, := \,
\sup_{x \xeof \Rd} \abs{ \, x^{\alpha} \partial^{\beta} \phi(x) \, } .
\end{equation}

\subparagraph*{}
Then $\phi$ is a Schwartz function
iff $p_{\alpha,\beta}(\phi) < \infty$
holds for all multi-indices $\alpha, \beta$.
The collection of mappings $p_{\alpha,\beta} \:\SRd \xto \R$
defines a countable family of seminorms,
which gives rise to a locally convex topology $\topT_{\mathscr{S}}$ on $\SRd$.
Moreover,
since the Schwartz space is complete {\wrt} $\topT_{\mathscr{S}}$,
it actually forms a \textit{{\FRECHET} space}.

\subparagraph*{}
Comparing the Schwartz space $\SRd$ with the function spaces
$\D(\Rd)$ and $\E(\Rd)$,
it is easy to see that there is the chain of inclusions
\[
\D(\Rd) \, \hookrightarrow \, \SRd \, \hookrightarrow \, \cal{E}(\Rd).
\]
Moreover,
by comparing the locally convex topologies of these spaces,
it turns out that these inclusions are actually continuous.

\subsubsection*{} % space of tempered distributions
Similar to the spaces $\D(\G)$,
we introduce the concept of \textit{tempered distribution}
as a continuous linear functional $T \:\SRd \xto \C$.
The space of tempered distributions is usually denoted $\TD$.
A tempered distribution generalizes the class of bounded and slow-growing
locally integrable functions.

\subparagraph*{}
The following criterion determines
under which conditions a linear functional $f$ on $\SRd$
forms a tempered distribution:
this is the case 
\iff
there is a constant $C > 0$ and integers $m,l \xeof \N$ such that
\[
\abs{ f(\phi) } \, \leq \,
C \cdot \sum_{\abs{\alpha} \leq l, \abs{\beta} \leq m} \, p_{\alpha,\beta}(\phi)
\qfa \phi \xeof \SRd\,,
\]
where the functions $p_{\alpha,\beta}$ are the seminorms
defined by \eqref{X_SRN_FAMILY_OF_SEMINORMS}.

\subsubsection*{} % comparison of distributional spaces
Comparing the known kinds of distributional spaces,
we obtain a chain of inclusions,
which is dual to that of the base spaces:
\[
\cal{E}'(\Rd) \, \xsubset \, \TD \, \xsubset \, \D'(\Rd).
\]
Endowing each of these spaces with the so-called \textit{weak*-topology},
these embeddings are continuous.
Hence a distribution having compact support can also be regarded
as a tempered distribution,
which in turn gives rise to a Schwartz distribution.

%------------------------------------------------------------------------------%
\section{Fourier Transforms in Distributional Spaces}                          %
%------------------------------------------------------------------------------%
In this section we shall elaborate that Fourier transformation can be
extended to the space $\TD$ of tempered distributions.
\subparagraph*{}
First notice 
that the Fourier transform $\hat{\phi}$
of a Schwartz function $\phi \xeof \SRd$ is well defined,
since $\SRd$ is a linear subspace of the convolution algebra $L^1(\Rd,+,\ast)$.
Moreover, 
Fourier transformation restricted to the Schwartz space $\SRd$
forms a linear mapping $\FT \:\SRd \xto \SRd$,
which sends each function $\phi \xeof \SRd$ to a further function
$\hat{\phi} \xeof \SRd$ according to the well-known formula
\begin{equation}\label{FOURIER_TRANSFORM_OF_SCHWARTZ_FUNCTION}
\hat{\phi}(x) \= \FT(\phi)(x) :=
\frac{1}{(2\pi)^{d/2}}
\int_{\,\Rd} e^{-i <\xi,x>} \phi(\xi) \, d \xi
\quad (x \xeof \Rd).
\end{equation}

\subparagraph*{}
The extraordinary meaning of Schwartz functions is founded by the fact that
the Fourier transformation has excellent properties on the space $\SRd$.
More precisely,
besides the well-known features of Fourier transformation in $L^1(\Rd,+,\ast)$,
the mapping $\FT$ is \textit{bijective} on $\SRd$.
Hence the \textit{inverse Fourier transform} exists for all $\phi \xeof \SRd$
and is described by the \textit{Fourier inversion formula}
\[
\check{\phi}(\xi) \= 
\FT^{-1}(\phi)(\xi) \eqdef
\frac{1}{(2\pi)^{d/2}}
\int_{\,\Rd} e^{i \scalar{x}{\xi}} \, \phi(x) \, dx
\quad \quad (\xi \xeof \Rd).
\]

\subparagraph*{}
We next extend Fourier transformation in $\SRd$ to the space $\TD$
of tempered distributions by transposition,
that is
\[
(\FT \disT) \, \phi \, := \, \disT(\FT \phi),
\quad \quad \; \phi \xeof \SRd, \; \disT \xeof \TD.
\]
This mapping is bijective in the space $\TD$
and even continuous {\wrt} the weak*-topology on $\TD$.
The inverse Fourier transform of a tempered distribution $\disT$
is represented by the formula
\[
\FT^{-1}(\disT)(\phi) \, := \, \disT(\FT^{-1} \phi)
\quad \quad (\phi \xeof \SRd).
\]
For example,
the Fourier transform $\FT(\delta_0)$ of the Dirac distribution
$\delta_0 \xeof \cal{S}'(\R)$ is the continuous linear functional
\[
\FT(\delta_0)(\phi) \, := \,
\delta_0(\FT(\phi)) \= \hat{\phi}(0) \=
\frac{1}{\sqrt{2\pi}} \, \int_{\,\R} \phi(x) \, dx,
\quad \phi \xeof \SR1\,,
\]
so $\FT(\delta_0)$ is nothing but the distribution induced by
the constant function $x \mapsto \frac{1}{\sqrt{2\pi}}$.

%------------------------------------------------------------------------------%
\section{Fundamental Solutions}                                                %
%------------------------------------------------------------------------------%
We now study the action of linear partial differential operators 
in distributional spaces.
First recall
that if $\psi \xeof \SF$ and $\disF \xeof \D'(\G)$,
then $\psi \disF$ defined by
\[
(\psi \disF) \phi \, := \, \disF (\psi \phi), \quad \quad \phi \xeof \TF
\]
also belongs to $\D'(\G)$. 
This implies that the linear differential operators $c_{\alpha} \, \PA$ 
with $c_{\alpha} \xeof \SF$ are acting on the space $\D'(\G)$,
since for $\disF \xeof \D'(\G)$
the mapping $c_{\alpha} \PA \disF$ is a distribution again.
By taking linear combinations of such operators,
we conclude that the action
of a linear partial differential operator
\[
\DO(x,u) \, := \,
\sum_{\abs{\alpha} \leq 2m} c_{\alpha}(x) \; \PA u(x)
\]
having smooth coefficients and acting on the space $\D'(\G)$ is well-defined,
which of course is especially true in case of $\G \eq \Rd$.

\subsubsection*{} % definition of fundamental solution
In the following we introduce the concept of \textit{fundamental solution}
for a linear partial differential operators $\DO$
having smooth coefficients $c_{\alpha}$.
It is nothing but a distribution $\Phi_{\DO} \xeof \D'(\Rd)$
that forms a solution of the differential equation
$\DO \Phi_{\DO} \eq \delta_0$,
where $\delta_0$ is the Dirac distribution {\wrt} the origin $x=0$.

\subparagraph*{}
It is useful to know a fundamental solution of $\DO$
as the following computation shows:
let $\disF \xeof \D'(\Rd)$ be any Schwartz distribution and consider the
inhomogeneous differential equation $\DO \, \disU \eq \disF$.
If convolution between $\Phi_{\DO}$ and $\disF$ exists,
then the equalities
\[
\DO(\Phi_{\DO} \ast \disF)  \=
(\DO \Phi_{\DO}) \ast \disF \, = \;
\delta_0 \ast \disF         \= \disF
\]
are valid,
so the convolutional product $\Phi_{\DO} \ast \disF$
solves the distributional differential equation $\DO \, \disU \eq \disF$. 
Notice that
there is a distributional solution $\disU$ of this equation
especially for all distributions $\disF$ having compact support.

\subparagraph*{}
The fundamental solution of a linear differential operator
need not be unique,
which can be showed by the following example:
let $\disU \xeof \D'(\Rd)$ be a solution of the homogeneous equation
$\DO \, \disU \eq 0$,
then the distribution $\Phi_{\DO} + \disU\,$
is a fundamental solution of $\DO$, too.

\subsubsection*{} % Malgrange-Ehrenpreis theorem
Regarding the existence of fundamental solutions for linear 
differential operators,
the \textsc{Malgrange-Ehrenpreis} \textit{theorem}
has been shown to be a milestone in PDE theory:

\begin{theorem}\label{MALGRANGE-EHRENPREIS_THEOREM}
Each linear differential operator $\DO \neq 0$
with constant coefficients
has a fundamental solution in the distributional space $\D'(\Rd)$.
\end{theorem}

%\input{PROOF_MALGRANGE-EHRENPREIS_THEOREM}

\subparagraph*{}
It should be mentioned 
that there are counterexamples which demonstrate
that this theorem is not valid for differential operators
whose coefficients are non-constant.

\subsubsection*{} % examples of fundamental solutions
In the following we describe the fundamental solutions
of three linear differential operators,
namely
the \textit{Laplace operator},
the \textit{heat operator} and
the \textit{wave operator}.

\begin{itemize}
\item[1.]
The fundamental solution $\Phi_{\Delta}^d$ of the Laplacian $\Delta$
depends on the number of spatial coordinates.
Since $\Delta$ is invariant {\wrt} translations and rotations,
it is nearby to conjecture $\Phi_{\Delta}^d(x) \eq \phi_d(\abs{x} )$.
Indeed,
the functions 
\[
\Phi_{\Delta}(x) \=
\begin{cases}[1.5]
\quad \frac{1}{2} \abs{x}                         & (d=1)      \\
\quad - \frac{1}{2 \pi} \ln \abs{x}               & (d=2)      \\
\quad \frac{1}{(d-2) \, \omega_d } \abs{x} ^{2-d} & (d \geq 3) \\
\end{cases}
\]
are solutions of the differential equations
$\Delta \Phi_{\Delta}^d \eq \delta_0$ for $d \xeof \N$,
where $\omega_d$ is the Lebesgue-Borel measure of
the $d$-dimensional unit sphere.
Each other fundamental solution of the Laplacian is of the form
$\Phi_{\Delta}^d + P_1$,
where $P_1$ is any polynomial of degree less than two.

\item[2.]
The fundamental solution of the \textit{heat operator} 
$\DO := \partial_t - \Delta u$ in $\Rd$ is given 
by the so-called \textit{heat kernel} function
\[
\Phi_d(x,t) \=
\begin{cases}[1.5]
\quad \frac{1}{(4 \pi t)^{d/2}} \cdot e^{- \frac{\abs{x} ^2}{4t}},
         & \; x \xeof \Rd, \; t > 0,   \\
\quad 0, & \; x \xeof \Rd, \; t \leq 0.\\
\end{cases}
\]
Each other fundamental solution of the heat operator
adopts the form $\Phi_d + c$ ($c \xeof \C$),
so the kernel of the heat operator 
purely consists of constant distributions.

\item[3.]
When investigating the fundamental solutions
of the hyperbolic \textit{wave operator} $\Box := \partial_{tt} - \Delta u$,
one has to distinguish between even and odd spatial dimensions.

\subparagraph*{}
Let $H$ be the \textit{Heaviside function},
defined by $H(t) := 0$ for $t<0$ and $H(t) := 1$ for $t \geq 0$.
Then we obtain for $d \leq 3$ 
\[
\Phi_{\Box}^d(x,t) \= \;
\begin{cases}[1.5]
\quad \frac{1}{2} H(t-\abs{x} )                                      & (d=1)\\
\quad \frac{1}{2 \pi} \, (t^2 - \abs{x} ^2)^{-1/2} \, H(t- \abs{x} ) & (d=2)\\
\quad \frac{1}{2 \pi} \, \delta(t^2 - \abs{x} ^2) \, H(t)            & (d=3).\\
\end{cases}
\]

\subparagraph*{}
It is a notable fact
that all the higher fundamental solutions of the wave operator
can be derived from the functions $\Phi_{\Box}^1$ and $\Phi_{\Box}^2$.
Using the variable $r := \abs{x} $,
the functions $\Phi_{\Box}^{d+2}$ ($d \xeof \N$) are recursively given by
\[
\Phi_{\Box}^{d+2}(t,r) \=
\frac{-1}{2\pi r} \cdot \frac{\partial}{\partial_r} \Phi_{\Box}^d(t,r),
\quad d \geq 1.
\]
To determine the support of a fundamental solution $\Phi_{\Box}^d$,
we first define for fixed $x_0 \xeof \Rd$ the \textit{light cone}
\[
C^+(x_0) \, := \, \{(t,x): t > 0, \abs{x-x_0} < t \}
\]
and obtain the support of the fundamental solution according to
\[
\supp{\Phi_{\Box}^d} \=
\begin{cases}[1.5]
\quad \partial C^+(0),   \quad     d \mbox{ odd  or } d \geq 3,\\
\quad \overline{C^+(0)}, \quad\;\; d \mbox{ even or } d=1.\\
\end{cases}
\]
\end{itemize}

\subsubsection*{} % hypoelliptic differential operators
A linear differential operator $\DO$ is called \textit{hypoelliptic},
if {\wrt} the linear differential equation $\DO u \eq f$
the condition $f \xeof \SF$ implies $u \xeof \SF$.
Restricting to differential operators with constant coefficients,
hypoellipticity can be characterized
by means of a feature of the fundamental solution $\Phi_{\DO_c}$:
$\DO_c$ is hypoelliptic
\iff
$\Phi_{\DO_c} \eq \disF_f$,
where $\disF_f$ is a regular distribution 
with generating function $f$ being indefinitely differentiable
in every $x \neq 0$.

\subparagraph*{}
We next introduce the concept of \textit{singular support} $\Sigma(\disF)$
for $\disF \xeof \D'(\Rd)$:
it is the smallest subset $S \xeof \Rd$
such that $\disF$ remains regular on $\Rd \setminus S$.
Then hypoellipticity of linear differential operators
with constant coefficients are characterized as follows:

\begin{theorem}\label{THEOREM_ON_HYPOELLIPTIC_OPERATORS}
A linear partial differential operator $\DO_c$ with constant coefficients
is hypoelliptic
\iff
$\,\Sigma(\DO_c \, \disF) \eq \Sigma(\disF)$ for all $\disF \xeof \D'(\Rd)$.
\end{theorem}

\subparagraph*{}
Knowing the fundamental solutions of the elliptic operator $\Delta$ and
the parabolic heat operator $\partial_t-\Delta$,
it turns out that both operators are hypoelliptic.
But theorem \ref{THEOREM_ON_HYPOELLIPTIC_OPERATORS} implies
that the hyperbolic wave operator $\Box \eq \partial_{tt}-\Delta$
is not hypoelliptic.

%------------------------------------------------------------------------------%
\section{Free Space Problems}                                                  %
%------------------------------------------------------------------------------%
To show the benefit of knowing the fundamental solution
of a linear partial differential operator with constant coefficients,
we consider free space problems ($\G \eq \Rd$)
related to the Laplacian, the heat and the wave operator.

\begin{itemize}
\item[1.]
The \textit{Poisson equation} in the entire space $\Rd$ is given by
\begin{equation}\label{CLASSICAL_POISSON_EQUATION}
\Delta \, u \= - 4\pi \,f, \quad f \xeof C(\Rd).
\end{equation}

\subsubsection*{} % Poisson equation
A function $u \:\Rd \xto \R$ is called a \textit{classical solution}
of \eqref{CLASSICAL_POISSON_EQUATION},
if it is twice continuously differentiable
and fulfills the differential equation in each $x \xeof \Rd$.
However,
$\disU \xeof \D'(\Rd)$ is said to be a \textit{distributional solution}
of the Poisson equation,
if
\begin{equation}\label{DISTRIBUTIONAL_POISSON_EQUATION}
\Delta \, \disU \= - 4\pi \,\disF, \quad \disF \xeof \D'(\Rd)
\end{equation}
holds in the sense of distributions.

\subparagraph*{}
Let $\Phi_{\Delta}$ be the fundamental solution of the Laplacian.
If the convolutional product $\Phi_{\Delta} \ast \disF$ exists,
which is especially true for each regular distribution $\disF_f$
with generating function $f \xeof C_c(\Rd)$,
then the solution of \eqref{DISTRIBUTIONAL_POISSON_EQUATION}
is obtained by convolution,
that is,
$\disU \eq \Phi_{\Delta} \ast \disF$.

\subparagraph*{}
For the physically interesting case $d=3$
and for $f \xeof C_c(\R^3)$ we get a classical solution $u_f$,
called the \textit{Newton potential} for the function $f$:
\[
u_f(x) \=
\int_{\,\R^3} \, \frac{f(y)}{\abs{x-y } }  \, dy, \quad x \xeof \R^3.
\]

\item[2.]
By the \textit{heat equation} equation in $\Rd$ is meant
the inital value problem
\begin{equation}\label{CLASSICAL_HEAT_EQUATION}
\begin{cases}[1.2]
\quad u_t(t,x) - \Delta u(t,x) = f(t,x) & (x \xeof \Rd, \, t > 0) \\
\quad u(0,x) \= 0,                      & (x \xeof \Rd). \\
\end{cases}
\end{equation}

\subsubsection*{} % heat equation
We first define the concept of distributional solution
for problem \eqref{CLASSICAL_HEAT_EQUATION}:
let $\disU_0 \xeof \D'(\Rd)$ and $\disF \xeof \D'(\R^{d+1})$ are
Schwartz distributions satisfying 
$\supp{\disF} \xsubset [0,\infty) \xtimes \Rd$,
then $\disU \xeof \D'(\R^{d+1})$ is called a \textit{distributional solution},
if $\supp{\disU} \xsubset [0,\infty) \xtimes \Rd$ and
\begin{equation}\label{DISTRIBUTIONAL_HEAT_EQUATION}
\disU_t \, - \Delta \, \disU \= \disF + \disU_0 \otimes \delta_0,
\end{equation}
where $\delta_0 \xeof \D'(\R)$ is the Dirac distribution.
Here the initial condition of the problem has been transformed
into an inhomogeneous part of the differential equation.

\subparagraph*{}
The existence of a solution for \eqref{DISTRIBUTIONAL_HEAT_EQUATION}
is ensured,
if $\disU_0 \xeof \E'(\Rd)$ and $\disF \xeof \E'(\R^{d+1})$.
In this case the solution $\disU \xeof \D'(\Rd)$ is given by
convolution of the fundamental equation $\Phi_d$ 
with the inhomogeneous parts of \eqref{DISTRIBUTIONAL_HEAT_EQUATION},
\[
\disU \=
\Phi_d \ast (\disF + \disU_0 \otimes \delta_0) \=
\Phi_d \ast \disF \, + \, \Phi_d \ast \disU_0.
\]

\subparagraph*{}
Given a classical solution $u$ of \eqref{CLASSICAL_HEAT_EQUATION},
the related Schwartz distribution $\disF_u$ is a solution
in the distributional sense.
Conversely,
let $\disF_u$ be regular and solve \eqref{DISTRIBUTIONAL_HEAT_EQUATION},
then $u$ forms a classical solution,
if some regularity conditions are valid:

\begin{itemize}[leftmargin=10mm]
\item[a.]
$\disU, \disU_0$ and $\disF$ are regular distributions
with associated functions $u, u_0, f \xeof \Lc$,

\item[b.]
$u_0$ is continuous on $\Rd$ and
$f$ is continuous on $[0,\infty) \xtimes \Rd$,

\item[c.]
$u$ is twice continuously differentiable on $[0,\infty) \xtimes \Rd$.
\end{itemize}

\item[3.]
The initial-value problem for the \textit{wave equation} in $\Rd$
is given by

\begin{equation}\label{CLASSICAL_WAVE_EQUATION}
\begin{cases}[1.2]
\quad u_{tt}(t,x) - \Delta u(t,x) \= f(t,x) & (x \xeof \Rd, \, t > 0) \\
\quad u(0,x) = u_0(x)                       & (x \xeof \Rd) \\
\quad u_t(0,x) = u_1(x)                     & (x \xeof \Rd. \\
\end{cases}
\end{equation}

\subsubsection*{} % wave equation
As for the heat equation,
we define for \eqref{CLASSICAL_WAVE_EQUATION}
the concept of distributional solution,
where again both initial value conditions are transformed
into inhomogeneous parts of the equation by means of tensor products.

\paragraph*{}
Let $\disU_0 \xeof \D'(\Rd)$, $\disU_1 \xeof \D'(\Rd)$
and $\disF \xeof \D'(\R^{d+1})$ be distributions,
then $\disU$ is called a distributional solution
of \eqref{CLASSICAL_WAVE_EQUATION},
if
\begin{equation}\label{DISTRIBUTIONAL_WAVE_EQUATION}
\disU_{tt} - \Delta \disU \=
\disF + \disU_0 \otimes \delta_0' \, + \, \disU_1 \otimes \delta_0 \,,
\end{equation}
where $\delta_0$ is the Dirac distribution for $t=0$ 
and $\delta_0'$ is the distributional derivative of $\delta_0$
{\wrt} the variable $t$.

\subparagraph*{}
Since the support of the function
$\disU_0 \, \otimes \, \delta_0' \, + \, \disU_1 \, \otimes \, \delta_0$
belongs to the set $\R^{d+1}$,
there is to each distribution $\disF \xeof \D'(\R^{d+1})$ with support
\[
\supp{\disF} \= \{\, (x,t) \xeof \Rd \xtimes \R: t \geq 0 \,\}
\]
a further distribution
$
\disU_F :=
\Phi_{\Box}^d \ast (\disF + \disU_0 \otimes \delta'_0
+ \disU_1 \otimes \delta_0)
\in \D'(\R^{d+1}),
$
which in fact solves the distributional differential equation
\eqref{DISTRIBUTIONAL_WAVE_EQUATION}.

\subparagraph*{}
Each classical solution of \eqref{CLASSICAL_WAVE_EQUATION}
is also a distributional solution.
However,
the reverse statement only holds under the conditions 

\begin{itemize}
\item[a.]
$\disU, \disU_0, \disU_1$ and $\disF$ are regular
and generated by ordinary functions $u, u_0, u_1, f \xeof \Lc$,

\item[b.]
$u_0 \xeof C^1(\Rd)$ and $u_1 \xeof C(\Rd)$,

\item[c.]
$f \xeof C([0,\infty) \xtimes \Rd)$ and

\item[d.]
$u \xeof C^2((0,\infty) \xtimes \Rd) \cap C([0,\infty) \xtimes \Rd)$.
\end{itemize}
\end{itemize}

%------------------------------------------------------------------------------%
\section{Applications of the Fourier Transform}                                %
%------------------------------------------------------------------------------%
The usability of Fourier transformation mainly relies on the fact
that elementary differential operators $\partial_k$,
applied to Schwartz functions $\phi$,
appear as multiplication operators to the Fourier transform $\FT(\phi)$
and via verce,
that is, 
we always have
\[
\partial_k \, \phi(x) \, \overset{\FT}{\longmapsto}
\, i \, \xi_k \cdot \hat{\phi}(\xi)\,,
\quad \quad
i \, x_k \phi(x) \, \overset{\FT}{\longmapsto}
\, \partial_k \, \hat{\phi}(\xi)\,,
\]
for $k=1,\ldots,d$,
as one can verify by partial integration.
For any linear differential operator $\DO_c$
with constant coefficients $c_{\alpha}$
and symbol $S_{\DO_c}$ we accordingly obtain 
\[
(\DO_c \, \phi)(x) \, \overset{\FT}{\longmapsto} \,
S_{\DO_c}(i \xi) \cdot \hat{\phi}(\xi).
\]
This is nothing but the fact
that all the differential operators $\DO_c$ are Fourier multipliers.
They operate on functions $\phi \xeof \SRd$ like multi\-plications
with the symbols $S_{\DO_c}$ on the Fourier trans\-form $\FT \phi$.
Hence a linear differential operator $\DO_c$ with constant coefficients
is conjugated to its symbol function $S_{\DO_c}$ in the sense of
\[
\DO_c \phi \=
(\FT^{-1} \circ S_{\DO_c} \circ \FT) \, \phi \=
\FT^{-1}(S_{\DO_c} \hat{\phi})
\qfa \phi \xeof \SRd.
\]

\subsubsection*{} % the general usage of the Fourier transform in PDE theory
We want incident
how Fourier transformation can be used
to solve linear differential equations in the whole space $\Rd$
and governed by the Laplacian.
The basic idea is to transform the differential equation
either into a purely algebraic equation
or into an ordinary differential equation.
Once these are solved,
inverse Fourier transformation is applied
to get the solution of the original problem.

\subparagraph*{}
Two examples for these methods will be presented.
In the first one,
the bijectivity of Fourier transformation on the Schwartz space
is used to solve a specific differential equation,
while in the second one the simplified computation of convolution
products after performing Fourier transformation is helpful
to find fundamental solutions of the underlying differential operators.

\begin{itemize}
\item[1.]
The equation $\DO_c \, u \eq f$ with $f \xeof \SRd$
and governed by a differential operator $\DO_c$ having constant coefficents
can be converted by applying Fourier transformation
into the algebraic equation
\begin{equation}\label{ALGEBRAIC_EQUATION_OF_LDE}
S_{\DO_c}(i \xi) \, \hat{u}(\xi) \= \hat{f}(\xi)
\quad \quad (\xi \xeof \Rd).
\end{equation}
Now,
if the symbol function $S_{\DO_c}$ of $\DO_c$ owns the property
\begin{equation}\label{SYMBOL_CONDITION}
S_{\DO_c}(i \xi) \neq 0
\qfa \xi = (\xi_1,\ldots,\xi_d) \xeof \Rd,
\end{equation}
equation \eqref{ALGEBRAIC_EQUATION_OF_LDE} can be resolved for $\hat{u}(\xi)$,
and the solution $u$ of $\DO_c \,u \eq f$
is obtained by applying inverse Fourier transformation:
\[
\hat{u}(\xi) = \hat{f}(\xi) \, / \, S_{\DO_c}(i \xi)
\]
and
\[
u(x) \= \frac{1}{(2 \pi)^{d/2}} \,
\int_{\,\Rd}
e^{i \xi x} \cdot S_{\DO_c}^{-1}(i \xi) \cdot \hat{f}(\xi) \, d \xi.
\]

\item[2.]
If \eqref{SYMBOL_CONDITION} holds,
the Fourier transform method can also be applied
to solve the distributional differential equation
$\DO_c \, \disU \eq \disF$ with $\disF \xeof \TD$.
Since $\DO_c$ is bijective on $\SRd$ and 
\[
P_{\DO_c}(i \xi) \= P_{\DO_c^*}(-\xi)
\qfa \xi \xeof \Rd
\]
holds,
the formally adjoint operator $\DO_c^*$
is bijective on $\SRd$ and so its transposed operator
$\DO_c \eq (\DO_c^*)^t$ is bijective on the dual space $\TD$.
Thus,
if \eqref{SYMBOL_CONDITION} is valid,
the equation $\DO_c \, \disU \eq \disF$
is uniquely solvable for each tempered distribution $\disF$.

\subparagraph*{}
Especially for the differential operator 
$\Delta{_\lambda} := -\Delta + \lambda$ and $\lambda > 0$ we obtain
\[
S_{\Delta_{\lambda}}(i \xi) = \abs{\xi} ^2 + \lambda \neq 0
\]
for $\xi \xeof \Rd$,
hence the solution $u$ of $\Delta_{\lambda} u \eq f$ with $f \xeof \SRd$ is
\[
u(x) \=
\frac{1}{(2 \pi)^{d/2}} \,
\int_{\,\Rd} e^{i \xi x}
\cdot \frac{\hat{f}(\xi)}{\abs{\xi} ^2 + \lambda} \, d \xi, \quad (x \xeof \Rd).
\]
\end{itemize}

\subsubsection*{} % partial Fourier transformation
Partial Fourier transformation is a method
where the operator $\FT$ is applied only to the spatial variables
of the partial differential equation.
The result is an ordinary differential equation
{\wrt} the time variable $t$.
Then solving this equation and applying inverse Fourier transformation
leads to the solution of the original problem.

\subsubsection*{} % examples of partial Fourier tranformation
In the following we sketch
how {partial Fourier transformation} is useful to solve
the initial-value problem of the heat and wave equation
in the entire space $\Rd$.

\begin{itemize}
\item[1.]
The initial-value problem of the heat equation is given by
\[
\begin{cases}[1.2]
\quad u_t - \Delta u = 0   & (t > 0)     \\
\quad u(0,x) \= f(x)       & (x \xeof \R). \\
\end{cases}
\]
Applying Fourier transformation to both sides of the differential equation
{\wrt} the spatial variables $x_1,\ldots,x_d$,
we obtain
\begin{equation}\label{PARTIAL_FOURIER_TRANSFORMATION_I}
\frac{\partial}{\partial t} \hat{u}(t,\xi) \, + \, \xi^2 \, \hat{u}(t,\xi)
\= 0,
\quad \quad \quad
\hat{u}(0,\xi) = \hat{f}(\xi),
\end{equation}
where $\xi = (\xi_1,\ldots,\xi_d)$ are the variables
in the \textit{phase space}.

\subparagraph*{}
Then multiplying both sides of the first equation
in \eqref{PARTIAL_FOURIER_TRANSFORMATION_I} by $e^{\abs{\xi} ^2 t}$ leads to
\[
\frac{\partial}{\partial t} \, e^{\abs{\xi} ^2 t} \, \hat{u}(t,\xi) \= 0
\]
with solution
$\hat{u}(t,\xi) \eq \phi(\xi) \, e^{- \abs{\xi} ^2 t}$.
The mapping $\xi \mapsto \phi(\xi)$ is
the integration constant for each fixed point $\xi$.
The second equation in \eqref{PARTIAL_FOURIER_TRANSFORMATION_I}
implies $\phi(\xi)=\hat{f}(\xi)$,
hence the Fourier trans\-formed solution $\hat{u}$ 
of the initial-value problem is nothing but
\begin{equation}\label{PARTIAL_FOURIER_TRANSFORMATION_II}
\hat{u}(t,\xi) \eq \hat{f}(\xi) \cdot e^{- \abs{\xi} ^2 t}.
\end{equation}

\subparagraph*{}
The solution of the original equation is obtained
by applying inverse Fourier transformation and the rule
$
\FT^{-1}(f \cdot g) \eq \FT^{-1}(f) \ast \FT^{-1}(g)
$
to each side of \eqref{PARTIAL_FOURIER_TRANSFORMATION_II}:
\[
u(t,x) \=
\FT^{-1}(\hat{f}(\xi) \cdot e^{- \abs{\xi} ^2 t}) \=
\FT^{-1}(\hat{f}(\xi)) * \FT^{-1}(e^{- \abs{\xi} ^2 t})
\]
\[
\= f(x) * \FT^{-1}(e^{- \abs{\xi} ^2 t})
\=
\frac{1}{\sqrt{4\pi t}}
\int_{\,\R} e^{-\abs{x-y} ^2/(4t)} \,f(y)\,dy.
\]
The function
\[
H_t(x,y) \, := \,
\frac{1}{\sqrt{4\pi t}} \, e^{- \abs{x-y} ^2/(4t)}
\=
\Phi_{\partial_t - \Delta}(t,x-y)
\]
is called a \textit{heat kernel},
since it constitutes the kernel of the integral operator
$\cal{H}_t:f \mapsto H_t \ast f$.

\item[2.]
The method also enables to solve the initial-value problem
for the wave equation
\[
\begin{cases}[1.2]
\quad u_{tt} - \, \Delta u \= 0  & (t > 0),    \\
\quad u(0,x) \= f(x)             & (x \xeof \Rd),\\
\quad u_t(0,x) \= g(x)           & (x \xeof \Rd).\\
\end{cases}
\]

\subparagraph*{}
Partial Fourier transformation let become this equation into
\[
\hat{u}_{tt}(t,\xi) \= \abs{\xi} ^2 \hat{u}(t,\xi),
\quad \quad
\hat{u}(0,\xi) \= \hat{f}(\xi),
\quad \quad
\hat{u_t}(0,\xi) \= \hat{g}(\xi),
\]
and the general solution of the ordinary differential equation
in phase space is 
\[
\hat{u}(t,\xi)
\=
\hat{f}(\xi) \, \cos( t \abs{\xi} )
\, + \,
\hat{g}(\xi) \: \frac{\sin(t \abs{\xi} )}{\abs{\xi} }.
\]

\subparagraph*{}
Unfortunately,
performing inverse Fourier transformation to $\hat{u}$
for any $n >1$ is not an easy task.
But for $n=1$ we obtain
\[
u_1(x,t) \=
\FT^{-1}\left(\hat{f}(\xi) \cdot \cos(\abs{\xi} ) \right)
\=
\frac{1}{2}\,[f(x+t) \, + \, f(x-t)],
\]
\[
u_2(x,t) \=
\FT^{-1}\left(\hat{g}(\xi) \cdot
\frac{\sin(\abs{\xi} )}{\abs{\xi} } \right)
\=
\frac{1}{2}\, \int_{-t}^{+t} g(x+s) \, ds,
\]
thus the solution of the wave equation is given by $u \eq u_1+u_2$,
known as the \textsc{d'Alembert} \textit{solution}
\[
u(t,x) \=
\frac{1}{2} \, [f(x+t) \, + \, f(x-t)]
\, + \,
\frac{1}{2t}\, \int_{-t}^{+t} g(x+s) \, ds.
\]
\end{itemize}

%------------------------------------------------------------------------------%
%                                 BIBLIOGRAPHY                                 %
%------------------------------------------------------------------------------%
%\newpage
%\thispagestyle{empty}
\vspace{10mm}

\titleformat{\section}
{\normalfont \bfseries \filcenter}
{\thesection.\;}{0mm}{}

\begin{thebibliography}{99}
\begin{small}
\bibitem{X_Yo}  K.\,Yosida:
\textit{{\FA}}, 5.\,Edition, Springer Verlag, 1976.
\end{small}
\end{thebibliography}

%------------------------------------------------------------------------------%
%                               END OF DOCUMENT                                %
%------------------------------------------------------------------------------%
\end{document}
